% Generated draft: requires mathematical and source review.
\documentclass[11pt]{article}
\usepackage[margin=25mm]{geometry}
\usepackage{fontspec}
\setmainfont{Latin Modern Roman}
\usepackage{amsmath,amssymb,mathtools,unicode-math}
\setmathfont{Latin Modern Math}
\usepackage{enumitem,needspace,xurl,hyperref}
\setlength{\parindent}{0pt}
\setlength{\parskip}{5pt}
\setlength{\emergencystretch}{3em}
\newcommand{\N}{\mathbb{N}}
\newcommand{\Z}{\mathbb{Z}}
\newcommand{\Q}{\mathbb{Q}}
\newcommand{\R}{\mathbb{R}}
\newcommand{\C}{\mathbb{C}}
\newcommand{\F}{\mathbb{F}}
\newcommand{\A}{\mathbb{A}}
\newcommand{\PP}{\mathbb{P}}
\newcommand{\E}{\mathbb{E}}
\newcommand{\eps}{\varepsilon}
\newcommand{\dd}{\,\mathrm{d}}
\begin{document}
\section*{UnsolvedMath ID: 30005161}
\textbf{Source problem code:} catalogue entry
\subsection{Definitions and mathematical statement}
For a measurable function \(f\) on \([0,1]^d\), the usual weak \(L^p\) quasi-norm, for \(0<p<\infty\), is
\[
\|f\|_{L^{p,\infty}([0,1]^d)}
:=\sup_{\lambda>0}\lambda\,\bigl|\{x\in[0,1]^d:|f(x)|>\lambda\}\bigr|^{1/p},
\]
where \(|\cdot|\) denotes Lebesgue measure. This convention supplies terminology; the record does not specify the intended range of \(p\).

The supplied statement introduces \(B_p(\Gamma)\) as the least correction factor in a weak \(L^p([0,1]^d)\) majorant inequality for frequencies \(\Gamma\subset[1,N]^d\). However, it does not give that inequality, its coefficient restrictions, or its comparison expression. Consequently, \(B_p(\Gamma)\) cannot be defined uniquely from the supplied record.

\textbf{Question, conditional on the missing definition.} For which families of frequency sets \(\Gamma_N\subset[1,N]^d\) does the intended correction factor satisfy
\[
\forall\varepsilon>0\quad\exists C_{\varepsilon}>0\quad
\forall N\text{ in the family's parameter range},\qquad
B_p(\Gamma_N)\le C_{\varepsilon}N^{\varepsilon}?
\]
Here \(C_{\varepsilon}\) must be independent of \(N\); its permitted dependence on \(p\), \(d\), and the family is unspecified. The family notation makes explicit the uniformity needed to interpret the supplied bound \(B_p(\Gamma)\ll N^{\varepsilon}\), but the source record does not identify the intended families or the range of \(N\).

Can the random-set methods attributed in the supplied curation to Mockenhaupt--Schlag, and referenced in the original excerpt as \([\mathrm{MS09}]\), be adapted to establish such estimates in the intended setting? The original excerpt asks only whether these methods can be adapted to ``this setting''; the necessary surrounding mathematical context is absent. A fully specified mathematical problem therefore requires recovery of the majorant inequality, its hypotheses, and the intended random-set construction.
\subsection{Short English statement}
The supplied question asks which frequency sets have a weak majorant correction factor growing more slowly than every positive power of \(N\), and whether Mockenhaupt--Schlag random-set methods can prove such bounds. The defining inequality and several parameter restrictions are missing, so this is a conditional statement requiring mathematical review.
\subsection{Sources}
\begin{itemize}
\item Local UnsolvedMath registry: supplied record with permanent ID \(30005161\), title \emph{Random Set Methods for Harmonic Analysis Estimates}, exact source problem code \texttt{catalogue entry}, in the Oberwolfach Reports Open Problems collection. The record supplies proposed year 2022 and no named proposer.
\item Supplied report citation: \emph{Real Analysis, Harmonic Analysis and Applications} (2023), Oberwolfach Reports, DOI \texttt{10.4171/owr/2022/31}, pp.~1743--1803. Supplied literature-source label: Oberwolfach Report 31/2022. \url{https://doi.org/10.4171/owr/2022/31}. The original excerpt includes the page indication 1781 and asks whether the methods of \([\mathrm{MS09}]\) can be adapted to the setting under discussion.
\item The supplied curation attributes the random-set methods to Mockenhaupt--Schlag. No full bibliographic entry for \([\mathrm{MS09}]\) is supplied, so its identity and precise methods require verification.
\item Literature checked at: \texttt{2026-08-22}. This is supplied dated curation, not an independently verified assessment. It labels the problem open and reports that no authoritative source located by that cutoff establishes the proposed adaptation. The background explicitly requests independent verification of time-sensitive status and cited claims.
\item The registry labels the statement \texttt{corrected\_verified} and reports that the official Oberwolfach report was checked and the notation repaired. These are supplied verification claims; no independent source consultation or mathematical verification is asserted here.
\end{itemize}
\end{document}
